% GENERATED FILE. Edit canonical lessons, then run npm run generate:books.
% canonical-source-hash: fnv1a64:61dde0cfb2278bad
% canonical-lessons: GE-C205-knochen, GE-C205-zahnarzt, GE-C205-sprechstunde, GE-C205-untersuchung, GE-C205-verletzung, GE-R205-first-pass-words-from-chapters-201-202, GE-R205-second-pass-words-from-chapters-203-205

\chapter{A bone, A dentist, Surgery hours, An examination, An injury}
\label{ch:ge-a-bone-a-dentist-surgery-hours-an-examination-an-injury}
\hlchaptermodality{205}

\begin{chapteropening}
\textbf{By the end of this chapter:} \emph{I can say a bone, a dentist, surgery hours, an examination and an injury.}

Complete the last lesson of chapter 205: i can say a bone, a dentist, surgery hours, an examination and an injury.
\end{chapteropening}

\section[der Knochen]{der Knochen — a bone}

\label{lesson:GE-C205-knochen}



\begin{quote}
Before the new one: say the German for plastic, then the German for the body.
\end{quote}

\subsection*{You'll want to know: der Knochen}
\textbf{der Knochen} — \textquotedblleft{}a bone\textquotedblright{}.

\begin{grammarlens}[title={What you've built}]
One more word for this chapter.
\end{grammarlens}

\subsection*{Your turn}
\begin{itemize}
\raggedright
  \item \emph{Say it:} \emph{der Knochen}
  \item \emph{Say it:} \emph{der Knochen}, once more
  \item \emph{Say it:} say \emph{der Knochen}
  \item \emph{Recall:} say the German for plastic, then the German for the body, then say \emph{der Knochen} again
\end{itemize}

\begin{tcolorbox}[breakable,colback=teal!4,colframe=teal!35!black,title={Before you move on}]
What does der Knochen mean? (\textquotedblleft{}A bone\textquotedblright{}.) Say it once more.
\end{tcolorbox}
\section[der Zahnarzt]{der Zahnarzt — a dentist}

\label{lesson:GE-C205-zahnarzt}



\begin{quote}
Before the new one: say the German for the body, then the German for a bone.
\end{quote}

\subsection*{You'll want to know: der Zahnarzt}
\textbf{der Zahnarzt} — \textquotedblleft{}a dentist\textquotedblright{}.

\begin{grammarlens}[title={What you've built}]
One more word for this chapter.
\end{grammarlens}

\subsection*{Your turn}
\begin{itemize}
\raggedright
  \item \emph{Say it:} \emph{der Zahnarzt}
  \item \emph{Say it:} \emph{der Zahnarzt}, once more
  \item \emph{Say it:} say \emph{der Zahnarzt}
  \item \emph{Recall:} say the German for the body, then the German for a bone, then say \emph{der Zahnarzt} again
\end{itemize}

\begin{tcolorbox}[breakable,colback=teal!4,colframe=teal!35!black,title={Before you move on}]
What does der Zahnarzt mean? (\textquotedblleft{}A dentist\textquotedblright{}.) Say it once more.
\end{tcolorbox}
\section[die Sprechstunde]{die Sprechstunde — surgery hours, office hours}

\label{lesson:GE-C205-sprechstunde}



\begin{quote}
Before the new one: say the German for a bone, then the German for a dentist.
\end{quote}

\subsection*{You'll want to know: die Sprechstunde}
\textbf{die Sprechstunde} — \textquotedblleft{}surgery hours, office hours\textquotedblright{}.

\begin{grammarlens}[title={What you've built}]
One more word for this chapter.
\end{grammarlens}

\subsection*{Your turn}
\begin{itemize}
\raggedright
  \item \emph{Say it:} \emph{die Sprechstunde}
  \item \emph{Say it:} \emph{die Sprechstunde}, once more
  \item \emph{Say it:} say \emph{die Sprechstunde}
  \item \emph{Recall:} say the German for a bone, then the German for a dentist, then say \emph{die Sprechstunde} again
\end{itemize}

\begin{tcolorbox}[breakable,colback=teal!4,colframe=teal!35!black,title={Before you move on}]
What does die Sprechstunde mean? (\textquotedblleft{}Surgery hours, office hours\textquotedblright{}.) Say it once more.
\end{tcolorbox}
\section[die Untersuchung]{die Untersuchung — an examination}

\label{lesson:GE-C205-untersuchung}



\begin{quote}
Before the new one: say the German for a dentist, then the German for surgery hours.
\end{quote}

\subsection*{You'll want to know: die Untersuchung}
\textbf{die Untersuchung} — \textquotedblleft{}an examination\textquotedblright{}.

\begin{grammarlens}[title={What you've built}]
One more word for this chapter.
\end{grammarlens}

\subsection*{Your turn}
\begin{itemize}
\raggedright
  \item \emph{Say it:} \emph{die Untersuchung}
  \item \emph{Say it:} \emph{die Untersuchung}, once more
  \item \emph{Say it:} say \emph{die Untersuchung}
  \item \emph{Recall:} say the German for a dentist, then the German for surgery hours, then say \emph{die Untersuchung} again
\end{itemize}

\begin{tcolorbox}[breakable,colback=teal!4,colframe=teal!35!black,title={Before you move on}]
What does die Untersuchung mean? (\textquotedblleft{}An examination\textquotedblright{}.) Say it once more.
\end{tcolorbox}
\section[die Verletzung]{die Verletzung — an injury}

\label{lesson:GE-C205-verletzung}



\begin{quote}
Before the new one: say the German for surgery hours, then the German for an examination.
\end{quote}

\subsection*{You'll want to know: die Verletzung}
\textbf{die Verletzung} — \textquotedblleft{}an injury\textquotedblright{}.

\begin{grammarlens}[title={What you've built}]
One more word for this chapter.
\end{grammarlens}

\subsection*{Your turn}
\begin{itemize}
\raggedright
  \item \emph{Say it:} \emph{die Verletzung}
  \item \emph{Say it:} \emph{die Verletzung}, once more
  \item \emph{Say it:} say \emph{die Verletzung}
  \item \emph{Recall:} say the German for surgery hours, then the German for an examination, then say \emph{die Verletzung} again
\end{itemize}

\begin{tcolorbox}[breakable,colback=teal!4,colframe=teal!35!black,title={Before you move on}]
What does die Verletzung mean? (\textquotedblleft{}An injury\textquotedblright{}.) Say it once more.
\end{tcolorbox}
\section[(dialogue)]{Words from chapters 201-202 — a first pass}

\label{lesson:GE-R205-first-pass-words-from-chapters-201-202}



\begin{quote}
Say the words for an examination and an injury.
\end{quote}

\subsection*{Your turn}
Say these aloud:

\begin{itemize}
\raggedright
  \item a tin, a pot, a frying pan, an oven, a dishwasher
  \item a vacuum cleaner, an iron, shampoo, a brush, a suit
\end{itemize}

\begin{tcolorbox}[breakable,colback=teal!4,colframe=teal!35!black,title={Before you move on}]
A tin? (\textbf{die Dose}.) A pot? (\textbf{der Topf}.) A frying pan? (\textbf{die Pfanne}.) An oven? (\textbf{der Ofen}.) A dishwasher? (\textbf{die Spülmaschine}.) A vacuum cleaner? (\textbf{der Staubsauger}.) An iron? (\textbf{das Bügeleisen}.) Shampoo? (\textbf{das Shampoo}.) A brush? (\textbf{die Bürste}.) A suit? (\textbf{der Anzug}.)
\end{tcolorbox}
\section[(dialogue)]{Words from chapters 203-205 — a second pass}

\label{lesson:GE-R205-second-pass-words-from-chapters-203-205}



\begin{quote}
Say the words for a purse and a size.
\end{quote}

\subsection*{Your turn}
Say these aloud:

\begin{itemize}
\raggedright
  \item a purse, a size, material, wool, silk
  \item leather, wood, metal, plastic, the body
  \item a bone, a dentist, surgery hours, an examination, an injury
\end{itemize}

\begin{tcolorbox}[breakable,colback=teal!4,colframe=teal!35!black,title={Before you move on}]
A purse? (\textbf{das Portemonnaie}.) A size? (\textbf{die Größe}.) Material? (\textbf{der Stoff}.) Wool? (\textbf{die Wolle}.) Silk? (\textbf{die Seide}.) Leather? (\textbf{das Leder}.) Wood? (\textbf{das Holz}.) Metal? (\textbf{das Metall}.) Plastic? (\textbf{der Kunststoff}.) The body? (\textbf{der Körper}.) A bone? (\textbf{der Knochen}.) A dentist? (\textbf{der Zahnarzt}.) Surgery hours? (\textbf{die Sprechstunde}.) An examination? (\textbf{die Untersuchung}.) An injury? (\textbf{die Verletzung}.)
\end{tcolorbox}
